\section{Контейнеризация и оркестрация с использованием Docker и Kubernetes}

\subsection{Контейнеризация в Docker}

Для воспроизводимой поставки и упрощения развёртывания сервис упакован в Docker-контейнер~\cite{docker_docs}.  
Контейнеризация реализована через многостадийную сборку, в которой компиляция выполняется на отдельной стадии сборки.  
В окружение сборки включаются необходимые инструменты (в частности, \texttt{protobuf-compiler} для генерации кода из \texttt{.proto}).  
В итоговый образ попадают только исполняемый файл и минимально необходимое окружение.  

Такой подход даёт практические преимущества:
\begin{itemize}
\item сборка Docker-образа снижает размер конечного образа;  
\item поставка сервиса становится независимой от окружения машины;  
\item запуск в Docker и Kubernetes упрощает эксплуатацию сервиса;  
\item конфигурирование через переменные окружения обеспечивает единый способ настройки.  
\end{itemize}

\subsection{Развёртывание в Kubernetes}

Для эксплуатации в оркестраторе подготовлен набор Kubernetes-манифестов для сервисного слоя~\cite{kubernetes_docs}.  
Типовой набор ресурсов включает:
\begin{itemize}
\item ресурс развёртывания (Deployment) управляет жизненным циклом pod-контейнеров и обеспечивает постепенное обновление (rolling update);  
\item ресурс Service публикует gRPC-эндпоинт в сети;  
\item ресурс HPA обеспечивает горизонтальное масштабирование по метрикам;  
\item ресурс PDB ограничивает деградацию доступности при обслуживании узлов;  
\item отдельный ресурс Job запускает нагрузочное тестирование.  
\end{itemize}

На рисунке~\ref{fig:k8s-if97-service} показана целостная схема
развёртывания сервисного слоя в Kubernetes.
Она объединяет внешний трафик,
Ingress/LoadBalancer,
ClusterIP-сервис, две реплики pod-контейнеров,
проверки состояния и отдельный Job для нагрузочного тестирования.
Такая структура отражает не только сетевой маршрут запроса,
но и эксплуатационные контуры, необходимые для обновления,
масштабирования и проверки состояния сервиса.

\begin{figure}[htbp]
\centering
\resizebox{\textwidth}{!}{
\begin{tikzpicture}[
    font=\small,
    >=Latex,
    box/.style={
        draw, rectangle, rounded corners=2pt, align=center,
        minimum width=3.5cm, minimum height=1.0cm, fill=white
    },
    servicebox/.style={
        draw, rectangle, rounded corners=2pt, align=center,
        minimum width=4.5cm, minimum height=1.1cm, fill=blue!5
    },
    podbox/.style={
        draw, rectangle, rounded corners=2pt, align=center,
        minimum width=6.8cm, minimum height=2.4cm, fill=gray!5
    },
    containerbox/.style={
        draw, rectangle, align=center,
        text width=6.2cm, minimum height=1.0cm, fill=white
    },
    group/.style={
        draw, dashed, rounded corners=3pt, inner sep=12pt
    },
    flow/.style={
        draw, -Latex, line width=0.9pt
    },
    probe/.style={
        draw, dashed, -Latex, line width=0.8pt
    }
]

% Внешние компоненты (подняли выше на 0.5см для запаса)
\node[box] (external) at (10, 14.5) {Внешний трафик};
\node[box, minimum width=4.0cm] (ingress) at (10, 13.0) {Ingress / LoadBalancer};

% Сервис
\node[servicebox] (svc) at (10, 8.5) {Service\\\texttt{if97-calculator-service}\\ClusterIP :50051};

% Job Pod
\node[podbox] (jobpod) at (2, 8.5) {};
\node[anchor=north, font=\small\bfseries, inner sep=4pt] at (jobpod.north) {Под задания: \texttt{grpc-loadtest}};
\node[containerbox] (jobctr) at ([yshift=-0.35cm]jobpod.center) {\texttt{grpc\_loadtest}\\target: \texttt{if97-calculator-service:50051}};

% Pod 1
\node[podbox] (pod1) at (6, 4.5) {};
\node[anchor=north, font=\small\bfseries, inner sep=4pt] at (pod1.north) {Под-контейнер: \texttt{if97-service}};
\node[containerbox] (ctr1) at ([yshift=-0.35cm]pod1.center) {\texttt{if97-calculator-service:1.1.0}\\gRPC :50051};

% Pod 2
\node[podbox] (pod2) at (14, 4.5) {};
\node[anchor=north, font=\small\bfseries, inner sep=4pt] at (pod2.north) {Под-контейнер: \texttt{if97-service}};
\node[containerbox] (ctr2) at ([yshift=-0.35cm]pod2.center) {\texttt{if97-calculator-service:1.1.0}\\gRPC :50051};

% Health Probes
\node[box, minimum width=4.0cm, minimum height=1.2cm] (probes) at (10, 1.0) {gRPC Health Probes\\(readiness / liveness)};

% === Линии связи с жесткими привязками (чтобы не "плыли") ===
\draw[flow] (external.south) -- (ingress.north);
\draw[flow] (ingress.south) -- (svc.north);
\draw[flow] (jobpod.east) -- (svc.west);

% Теперь стрелки выходят из одной точки (нижний центр сервиса)
\draw[flow] (svc.south) -- (pod1.north);
\draw[flow] (svc.south) -- (pod2.north);

\draw[probe] (probes.north west) -- (pod1.south);
\draw[probe] (probes.north east) -- (pod2.south);


% === Группирующие рамки ===

% 1. Deployment
\node[group, fit=(pod1)(pod2)] (deploygrp) {};
\node[fill=white, inner sep=4pt, font=\small] at (deploygrp.north) {Развёртывание (Deployment) \texttt{if97-calculator-service} (\texttt{replicas: 2}, HPA 2--10)};

% 2. Namespace
\node (dummy_ns) at (2, 10.3) {}; 
\node[group, fit=(svc)(deploygrp)(jobpod)(dummy_ns)] (nsgrp) {};
\node[anchor=north west, font=\large\bfseries] at ([xshift=8pt, yshift=-8pt]nsgrp.north west) {Namespace};

% 3. Kubernetes Node (Верхнюю границу зафиксировали ниже, чтобы не пересекалась с Ingress)
\node (dummy_node) at (2, 11.2) {};
\node[group, fit=(nsgrp)(probes)(dummy_node)] (nodegrp) {};
\node[anchor=north west, font=\large\bfseries] at ([xshift=8pt, yshift=-8pt]nodegrp.north west) {Узел Kubernetes};

\end{tikzpicture}
}
\caption{Схема развёртывания сервисного слоя \texttt{if97\_calculator\_service} в Kubernetes: внешний трафик, сетевой вход, сервис, реплики pod-контейнеров, проверки состояния и нагрузочное задание (Job)}
\label{fig:k8s-if97-service}
\end{figure}

При развёртывании важны два эксплуатационных свойства.  
Первое свойство --- масштабируемость: при увеличении нагрузки можно увеличивать число реплик сервиса и настраивать лимиты вычислительного пула на уровне конфигурации.  
Второе свойство --- корректная остановка: при обновлениях и перемещениях pod-контейнеров сервис должен завершать активные вычисления предсказуемо, используя проверку работоспособности и режим штатного опустошения очереди.  
